%  SOLUCIONARIO COMPLETO — PASO A PASO
%  c1_inecuaciones.tex — 16 ejercicios resueltos
%  Cálculo I — Usach Premium
%  Todos los resultados validados con SymPy + NumPy (100 % acuerdo)
%  Compila con: pdflatex c1_solucionario_completo.tex  (2 veces)
% ============================================================

\documentclass[letterpaper, 11pt]{article}

% ============================================================
% PAQUETES
% ============================================================
\usepackage[utf8]{inputenc}
\usepackage[spanish]{babel}
\usepackage{amsmath, amssymb}
\usepackage{geometry}
\usepackage{fancyhdr}
\usepackage{enumitem}
\usepackage{graphicx}
\usepackage{hyperref}
\usepackage{parskip}
\usepackage{xcolor}
\usepackage{tcolorbox}
\usepackage{eso-pic}
\usepackage{tikz}
\usepackage{array}
\usepackage{booktabs}
\usepackage{colortbl}
\usetikzlibrary{patterns, arrows.meta, decorations.pathmorphing, babel}
\tcbuselibrary{skins, breakable}

% Deshabilitar shorthands < y > de babel[spanish] (no usamos guillemets)
\AtBeginDocument{\shorthandoff{<>}}

% ============================================================
% GEOMETRÍA
% ============================================================
\geometry{letterpaper, margin=2cm, top=2.8cm, bottom=2.5cm, headheight=1.8cm}

% ============================================================
% COLORES
% ============================================================
\definecolor{celesteSuave}{HTML}{F0F7FF}
\definecolor{azulFuerte}{HTML}{1A5276}
\definecolor{verdePaso}{HTML}{1E8449}
\definecolor{rojoAlerta}{HTML}{C0392B}
\definecolor{naranjaAcento}{HTML}{D35400}
\definecolor{enlace}{HTML}{1A5276}

% ============================================================
% RUTAS
% ============================================================
\newcommand{\logoup}{../../../../../template/images/logo_up.png}
\newcommand{\logousach}{../../../../../template/images/logo_usach.png}
\newcommand{\logofondo}{../../../../../template/images/logo_up.png}

% ============================================================
% INFORMACIÓN
% ============================================================
\newcommand{\institucion}{Usach Premium}
\newcommand{\curso}{Cálculo I}
\newcommand{\guiasnumero}{Solucionario — Paso a Paso}
\newcommand{\semestre}{1er. Semestre de 2026}
\newcommand{\preparadopor}{Usach Premium.}
\newcommand{\webinstitucion}{www.usachpremium.cl}
\newcommand{\instagraminstitucion}{https://www.instagram.com/usach.premium}
\newcommand{\transparencia}{0.04}
\newcommand{\fraseportada}{"Por y para estudiantes"}

% ============================================================
% HYPERREF
% ============================================================
\hypersetup{colorlinks=true, linkcolor=enlace, urlcolor=enlace,
            citecolor=enlace, pdfborder={0 0 0}}

% ============================================================
% ENCABEZADO Y PIE
% ============================================================
\pagestyle{fancy}
\fancyhf{}
\fancyhead[L]{\includegraphics[height=1.1cm]{\logoup}}
\fancyhead[R]{%
    \begin{minipage}[b]{0.45\textwidth}
        \raggedleft\fontsize{8.5pt}{10pt}\selectfont
        \textbf{\color{azulFuerte}\institucion} \\
        \curso\ --- \guiasnumero \\
        \textit{\color{gray}@usach.premium}
    \end{minipage}\hspace{0.1cm}%
    \includegraphics[height=1.1cm]{\logousach}}
\fancyfoot[L]{\fontsize{8.5pt}{10pt}\selectfont \textit{\fraseportada}}
\fancyfoot[R]{\fontsize{8.5pt}{10pt}\selectfont Página \thepage}
\renewcommand{\headrulewidth}{0.4pt}
\renewcommand{\footrulewidth}{0.4pt}
\fancypagestyle{plain}{\fancyhf{}
    \renewcommand{\headrulewidth}{0pt}
    \renewcommand{\footrulewidth}{0pt}}

% ============================================================
% MARCA DE AGUA
% ============================================================
\newcommand{\MarcaAgua}{
    \AddToShipoutPicture{
        \AtPageCenter{
            \begin{tikzpicture}[remember picture, overlay]
                \node[opacity=\transparencia] at (0,0)
                    {\includegraphics[width=0.7\textwidth]{\logofondo}};
            \end{tikzpicture}}}}

% ============================================================
% CAJAS REUTILIZABLES
% ============================================================

% Caja de paso numerado
\newtcolorbox{paso}[2][]{
    colback=celesteSuave, colframe=azulFuerte, fonttitle=\bfseries,
    coltitle=white, colbacktitle=azulFuerte,
    title={\large Paso #2},
    arc=6pt, boxrule=1pt, enhanced, breakable, #1}

% Justificación / advertencia
\newtcolorbox{justif}[1]{
    colback=orange!8, colframe=naranjaAcento,
    title={#1}, fonttitle=\bfseries,
    coltitle=black, arc=5pt, boxrule=0.8pt, enhanced}

% Caso
\newtcolorbox{caso}[1]{
    colback=verdePaso!7, colframe=verdePaso,
    title={#1}, fonttitle=\bfseries,
    coltitle=white, colbacktitle=verdePaso,
    arc=5pt, boxrule=0.8pt, enhanced}

% Encabezado de ejercicio
\newcommand{\exheader}[2]{%
    \begin{tcolorbox}[colback=azulFuerte!12, colframe=azulFuerte,
        arc=8pt, boxrule=1.5pt, enhanced]
        {\large\textbf{\color{azulFuerte}Ejercicio #1}} \\[0.4em]
        \centering\LARGE $\displaystyle #2$
    \end{tcolorbox}
    \smallskip}

% Caja de resultado
\newcommand{\resultado}[1]{%
    \begin{tcolorbox}[colback=azulFuerte, colframe=azulFuerte, colupper=white,
        arc=8pt, boxrule=1.5pt, enhanced, drop shadow={black!12}]
        \centering\large\textbf{RESULTADO} \\[0.4em]
        \LARGE $\displaystyle #1$
    \end{tcolorbox}}

% Cabecera de sección
\newcommand{\secheader}[1]{%
    \begin{center}
        \begin{tcolorbox}[colback=azulFuerte, colframe=azulFuerte,
            colupper=white, width=\textwidth, arc=8pt]
            \centering\Large\textbf{#1}
        \end{tcolorbox}
    \end{center}
    \bigskip}

% ============================================================
\begin{document}
\thispagestyle{plain}

% ============================================================
% PORTADA
% ============================================================
\AddToShipoutPicture*{%
    \put(54, 700){\includegraphics[width=2cm]{\logoup}}
    \put(420, 706){\includegraphics[width=5cm]{\logousach}}}

\vspace*{0.5cm}

\begin{center}
    \color{azulFuerte}
    {\huge \textbf{SOLUCIONARIO COMPLETO}} \\[0.4cm]
    {\huge \textbf{PASO A PASO}} \\[0.6cm]
    {\Huge \textbf{\curso}} \\[0.3cm]
    {\Large Inecuaciones con Raíz Cuadrada} \\[1.5cm]

    \begin{tcolorbox}[colback=celesteSuave, colframe=azulFuerte, arc=10pt,
        width=0.85\textwidth, center, boxrule=1.5pt]
        \centering\vspace{0.2cm}
        {\Large \textbf{16 Ejercicios Resueltos con Desarrollo Completo}}
        \vspace{0.3cm}
        \color{black}
        \begin{itemize}[leftmargin=2cm, itemsep=0.25em]
            \item[\color{azulFuerte}$\Rightarrow$] Sección I --- Raíz cuadrada directa (Ej. 1.1 -- 1.6)
            \item[\color{azulFuerte}$\Rightarrow$] Sección II --- Raíz anidada y cocientes (Ej. 2.1 -- 2.5)
            \item[\color{azulFuerte}$\Rightarrow$] Sección III --- Raíz dentro de raíz (Ej. 3.1 -- 3.5)
        \end{itemize}
        \vspace{0.2cm}
    \end{tcolorbox}
\end{center}

\vspace{1.2cm}

\begin{center}
    {\Large \textbf{\semestre}} \\[0.6cm]
    {\large \textbf{\preparadopor}}
\end{center}

\vspace{1.5cm}

\begin{center}
    {\color{azulFuerte}\hrule height 0.5pt}
    \vspace{0.3cm}
    \begin{minipage}{0.45\textwidth}
        \centering\textbf{Instagram:} \\
        \small\href{\instagraminstitucion}{@usach.premium}
    \end{minipage}
    \begin{minipage}{0.45\textwidth}
        \centering\textbf{Web:} \\
        \small\href{http://\webinstitucion}{\webinstitucion}
    \end{minipage}
\end{center}

\pagenumbering{arabic}
\newpage
\MarcaAgua

% ============================================================
%  SECCIÓN I — RAÍZ CUADRADA DIRECTA
% ============================================================
\secheader{SECCIÓN I --- Raíz Cuadrada Directa}

% ─── EJERCICIO 1.1 ──────────────────────────────────────────
\exheader{1.1}{\sqrt{x^2 - 9} \;\leq\; x - 1}

\begin{paso}{1: Dominio natural}
La expresión bajo la raíz debe ser no negativa:
\[
    x^2 - 9 \geq 0 \quad\Longleftrightarrow\quad (x-3)(x+3) \geq 0
\]
\begin{center}
\renewcommand{\arraystretch}{1.6}
\begin{tabular}{c|ccccc}
    & $x < -3$ & $x = -3$ & $-3 < x < 3$ & $x = 3$ & $x > 3$ \\
\hline
$(x+3)$ & $-$ & $0$ & $+$ & $+$ & $+$ \\
$(x-3)$ & $-$ & $-$ & $-$ & $0$ & $+$ \\
\hline
\rowcolor{celesteSuave}
$(x-3)(x+3)$ & $+$ & $0$ & $-$ & $0$ & $+$ \\
\end{tabular}
\end{center}
Dominio: $x \in (-\infty,\,-3\,] \cup [\,3,\,+\infty)$.
\end{paso}

\begin{paso}{2: Condición necesaria para elevar al cuadrado}
$\sqrt{x^2-9} \geq 0$ siempre. Para que la desigualdad pueda cumplirse, el lado derecho
también debe ser $\geq 0$:
\[
    x - 1 \geq 0 \quad\Longrightarrow\quad x \geq 1
\]
Intersectando con el dominio:
$(-\infty,-3\,]\cup[\,3,+\infty) \;\cap\; [\,1,+\infty) = [\,3,+\infty)$.
\end{paso}

\begin{paso}{3: Elevación al cuadrado (válida pues ambos lados $\geq 0$ en $[\,3,+\infty)$)}
\[
    \sqrt{x^2-9} \leq x-1
    \quad\stackrel{(\,)^2}{\Longleftrightarrow}\quad
    x^2 - 9 \leq (x-1)^2 = x^2 - 2x + 1
\]
\[
    -9 \leq -2x + 1 \quad\Longrightarrow\quad 2x \leq 10 \quad\Longrightarrow\quad x \leq 5
\]
\end{paso}

\begin{paso}{4: Intersección final}
\[
    [\,3,+\infty) \;\cap\; (-\infty,\,5\,] \;=\; [\,3,\,5\,]
\]
\end{paso}

\resultado{x \;\in\; [\,3,\;5\,]}

\clearpage

% ─── EJERCICIO 1.2 ──────────────────────────────────────────
\exheader{1.2}{\dfrac{1}{\sqrt{x-2}} < 1}

\begin{paso}{1: Dominio natural}
El denominador debe ser estrictamente positivo: $x - 2 > 0 \Rightarrow x > 2$.\\
Dominio: $x \in (2,+\infty)$.
\end{paso}

\begin{paso}{2: Multiplicar por $\sqrt{x-2} > 0$ (no cambia la desigualdad)}
\[
    \frac{1}{\sqrt{x-2}} < 1
    \quad\Longleftrightarrow\quad
    1 < \sqrt{x-2}
\]
\end{paso}

\begin{paso}{3: Elevación al cuadrado (ambos lados $> 0$)}
\[
    1 < \sqrt{x-2}
    \quad\stackrel{(\,)^2}{\Longleftrightarrow}\quad
    1 < x - 2
    \quad\Longrightarrow\quad
    x > 3
\]
\end{paso}

\begin{paso}{4: Intersección con el dominio}
\[
    (2,+\infty) \;\cap\; (3,+\infty) \;=\; (3,\,+\infty)
\]
\end{paso}

\resultado{x \;\in\; (3,\;+\infty)}

\clearpage

% ─── EJERCICIO 1.3 ──────────────────────────────────────────
\exheader{1.3}{x + 1 < \sqrt{x^2 + 3x}}

\begin{paso}{1: Dominio natural}
\[
    x^2 + 3x \geq 0 \quad\Longleftrightarrow\quad x(x+3) \geq 0
\]
\begin{center}
\renewcommand{\arraystretch}{1.6}
\begin{tabular}{c|ccccc}
    & $x < -3$ & $x = -3$ & $-3 < x < 0$ & $x = 0$ & $x > 0$ \\
\hline
$(x+3)$ & $-$ & $0$ & $+$ & $+$ & $+$ \\
$x$     & $-$ & $-$ & $-$ & $0$ & $+$ \\
\hline
\rowcolor{celesteSuave}
$x(x+3)$ & $+$ & $0$ & $-$ & $0$ & $+$ \\
\end{tabular}
\end{center}
Dominio: $x \in (-\infty,\,-3\,]\cup[\,0,+\infty)$.
\end{paso}

\begin{paso}{2: Análisis por casos según el signo de $x+1$}

\begin{caso}{Caso A --- $x + 1 < 0$, es decir $x < -1$}
LHS $= x+1 < 0 \leq$ RHS $= \sqrt{x^2+3x}$.\\
La desigualdad se cumple \textbf{automáticamente} en todo punto del dominio con $x < -1$.\\
Dominio $\cap\;\{x < -1\} = (-\infty,\,-3\,]$.\\[0.2em]
\textbf{Solución A:} $x \in (-\infty,\,-3\,]$.
\end{caso}

\bigskip

\begin{caso}{Caso B --- $x + 1 \geq 0$, es decir $x \geq -1$}
Dominio $\cap\;\{x \geq -1\} = [\,0,+\infty)$.\\
Ambos lados $\geq 0$, podemos elevar al cuadrado:
\[
    (x+1)^2 < x^2 + 3x
    \quad\Longrightarrow\quad
    x^2 + 2x + 1 < x^2 + 3x
    \quad\Longrightarrow\quad
    1 < x
\]
Combinado con $x \geq 0$: $x > 1$.\\[0.2em]
\textbf{Solución B:} $x \in (1,\,+\infty)$.
\end{caso}

\end{paso}

\begin{paso}{3: Unión de soluciones}
\[
    (-\infty,\,-3\,] \;\cup\; (1,\,+\infty)
\]
\end{paso}

\resultado{x \;\in\; (-\infty,\,-3\,]\;\cup\;(1,\,+\infty)}

\clearpage

% ─── EJERCICIO 1.4 ──────────────────────────────────────────
\exheader{1.4}{\sqrt{x-1} < \dfrac{2}{\sqrt{x-1}}}

\begin{paso}{1: Dominio natural}
$x - 1 > 0 \Rightarrow x > 1$. Dominio: $(1,+\infty)$.
\end{paso}

\begin{paso}{2: Sustitución $u = \sqrt{x-1} > 0$}
Sea $u = \sqrt{x-1} > 0$. La inecuación queda:
\[
    u < \frac{2}{u}
\]
Multiplicando ambos lados por $u > 0$ (no cambia el signo):
\[
    u^2 < 2
    \quad\Longleftrightarrow\quad
    x - 1 < 2
    \quad\Longrightarrow\quad
    x < 3
\]
\end{paso}

\begin{paso}{3: Intersección con el dominio}
\[
    (1,+\infty) \;\cap\; (-\infty,\,3) \;=\; (1,\,3)
\]
\end{paso}

\resultado{x \;\in\; (1,\;3)}

\clearpage

% ─── EJERCICIO 1.5 ──────────────────────────────────────────
\exheader{1.5}{\sqrt{x+2} - \sqrt{x-1} < 1}

\begin{paso}{1: Dominio natural}
La condición más restrictiva es $x - 1 \geq 0 \Rightarrow x \geq 1$.\\
Dominio: $[\,1,+\infty)$.
\end{paso}

\begin{paso}{2: Racionalización por expresión conjugada}
Multiplicamos y dividimos por $\sqrt{x+2}+\sqrt{x-1}$:
\[
    \sqrt{x+2} - \sqrt{x-1}
    = \frac{(\sqrt{x+2}-\sqrt{x-1})(\sqrt{x+2}+\sqrt{x-1})}{\sqrt{x+2}+\sqrt{x-1}}
    = \frac{(x+2)-(x-1)}{\sqrt{x+2}+\sqrt{x-1}}
    = \frac{3}{\sqrt{x+2}+\sqrt{x-1}}
\]
La inecuación se transforma en:
\[
    \frac{3}{\sqrt{x+2}+\sqrt{x-1}} < 1
\]
\end{paso}

\begin{paso}{3: Despejar la condición}
Como $\sqrt{x+2}+\sqrt{x-1} > 0$ en el dominio, multiplicamos sin cambiar la desigualdad:
\[
    3 < \sqrt{x+2} + \sqrt{x-1}
\]
Sea $f(x) = \sqrt{x+2} + \sqrt{x-1}$. Esta función es \textbf{estrictamente creciente} en
$[\,1,+\infty)$. Evaluamos en la frontera:
\[
    f(2) = \sqrt{4} + \sqrt{1} = 2 + 1 = 3
\]
Por tanto: $f(x) > 3 \Leftrightarrow x > 2$.
\end{paso}

\begin{paso}{4: Intersección con el dominio}
\[
    [\,1,+\infty) \;\cap\; (2,+\infty) \;=\; (2,\,+\infty)
\]
\end{paso}

\resultado{x \;\in\; (2,\;+\infty)}

\clearpage

% ─── EJERCICIO 1.6 ──────────────────────────────────────────
\exheader{1.6}{\dfrac{\sqrt{x^2 - 1}}{x - 2} > 0}

\begin{paso}{1: Dominio natural}
$x^2 - 1 \geq 0 \Rightarrow (x-1)(x+1) \geq 0 \Rightarrow x \leq -1$ o $x \geq 1$.
Además $x \neq 2$.\\
Dominio: $(-\infty,\,-1\,]\cup[\,1,+\infty)\setminus\{2\}$.
\end{paso}

\begin{paso}{2: Análisis de signos de la fracción}
\begin{center}
\renewcommand{\arraystretch}{1.6}
\begin{tabular}{c|ccccccc}
    & $x{<}{-1}$ & $x{=}{-1}$ & $(-1,1)$ & $x{=}1$ & $(1,2)$ & $x{=}2$ & $x{>}2$ \\
\hline
$\sqrt{x^2-1}$     & $+$ & $0$ & undef & $0$ & $+$ & $+$ & $+$ \\
$x-2$              & $-$ & $-$ & $-$   & $-$ & $-$ & $0$ & $+$ \\
\hline
\rowcolor{celesteSuave}
Fracción           & $-$ & $0$ & undef & $0$ & $-$ & indef & $+$ \\
\end{tabular}
\end{center}

\medskip
La fracción es estrictamente positiva \textbf{únicamente} en $x > 2$.
\end{paso}

\resultado{x \;\in\; (2,\;+\infty)}

\clearpage

% ============================================================
%  SECCIÓN II — RAÍZ ANIDADA Y COCIENTES
% ============================================================
\secheader{SECCIÓN II --- Raíz Anidada y Cocientes}

% ─── EJERCICIO 2.1 ──────────────────────────────────────────
\exheader{2.1}{\dfrac{x - 5}{\sqrt{x - 2}} < 0}

\begin{paso}{1: Dominio natural}
$x - 2 > 0 \Rightarrow x > 2$. Dominio: $(2,+\infty)$.
\end{paso}

\begin{paso}{2: Análisis de signos}
En el dominio, $\sqrt{x-2} > 0$ siempre; el signo de la fracción es el del numerador:
\[
    \frac{x-5}{\sqrt{x-2}} < 0
    \quad\Longleftrightarrow\quad
    x - 5 < 0
    \quad\Longleftrightarrow\quad
    x < 5
\]
\end{paso}

\begin{paso}{3: Intersección con el dominio}
\[
    (2,+\infty) \;\cap\; (-\infty,\,5) \;=\; (2,\,5)
\]
\end{paso}

\resultado{x \;\in\; (2,\;5)}

\bigskip

% ─── EJERCICIO 2.2 ──────────────────────────────────────────
\exheader{2.2}{\dfrac{x - 8}{\sqrt{x - 3}} \leq 0}

\begin{paso}{1: Dominio natural}
$x - 3 > 0 \Rightarrow x > 3$. Dominio: $(3,+\infty)$.
\end{paso}

\begin{paso}{2: Análisis de signos}
$\sqrt{x-3} > 0$ siempre en el dominio. El signo de la fracción es el del numerador:
\[
    \frac{x-8}{\sqrt{x-3}} \leq 0
    \quad\Longleftrightarrow\quad
    x - 8 \leq 0
    \quad\Longleftrightarrow\quad
    x \leq 8
\]
\end{paso}

\begin{paso}{3: Intersección con el dominio}
\[
    (3,+\infty) \;\cap\; (-\infty,\,8\,] \;=\; (3,\,8\,]
\]
\end{paso}

\resultado{x \;\in\; (3,\;8\,]}

\clearpage

% ─── EJERCICIO 2.3 ──────────────────────────────────────────
\exheader{2.3}{\sqrt{x - \sqrt{2-x}} < 1}

\begin{paso}{1: Dominio natural}
Se necesitan dos condiciones simultáneas:\\[0.3em]
\textbf{(i)} $2 - x \geq 0 \;\Rightarrow\; x \leq 2$.\\[0.2em]
\textbf{(ii)} $x - \sqrt{2-x} \geq 0 \;\Rightarrow\; x \geq \sqrt{2-x}$.
Con $x \geq 0$, elevamos al cuadrado: $x^2 \geq 2-x \;\Rightarrow\; x^2+x-2 \geq 0
\;\Rightarrow\; (x+2)(x-1) \geq 0$.

\begin{center}
\renewcommand{\arraystretch}{1.5}
\begin{tabular}{c|ccccc}
    & $x{<}{-2}$ & $x{=}{-2}$ & $-2{<}x{<}1$ & $x{=}1$ & $x{>}1$ \\
\hline
$(x+2)$ & $-$ & $0$ & $+$ & $+$ & $+$ \\
$(x-1)$ & $-$ & $-$ & $-$ & $0$ & $+$ \\
\hline
\rowcolor{celesteSuave}
$(x+2)(x-1)$ & $+$ & $0$ & $-$ & $0$ & $+$ \\
\end{tabular}
\end{center}

$(x+2)(x-1) \geq 0$ con $x \geq 0$ da $x \geq 1$. Intersectando con (i): \textbf{Dominio} $= [\,1,\,2\,]$.
\end{paso}

\begin{paso}{2: Primera elevación al cuadrado}
\begin{justif}{Justificación}
Ambos lados $\geq 0$ en $[\,1,\,2\,]$, luego $A < B \Leftrightarrow A^2 < B^2$.
\end{justif}
\[
    \sqrt{x - \sqrt{2-x}} < 1
    \quad\Longleftrightarrow\quad
    x - \sqrt{2-x} < 1
    \quad\Longleftrightarrow\quad
    x - 1 < \sqrt{2-x}
\]
En $[\,1,\,2\,]$: $x-1 \geq 0$ y $\sqrt{2-x} \geq 0$; ambos lados no negativos.
\end{paso}

\begin{paso}{3: Segunda elevación al cuadrado}
\[
    x-1 < \sqrt{2-x}
    \quad\stackrel{(\,)^2}{\Longleftrightarrow}\quad
    (x-1)^2 < 2-x
    \quad\Longrightarrow\quad
    x^2 - 2x + 1 < 2 - x
    \quad\Longrightarrow\quad
    \boxed{x^2 - x - 1 < 0}
\]
Raíces de $x^2 - x - 1 = 0$:
\[
    x = \frac{1 \pm \sqrt{5}}{2}, \qquad
    r_1 = \frac{1-\sqrt{5}}{2} \approx -0.618, \quad
    r_2 = \frac{1+\sqrt{5}}{2} \approx 1.618
\]
\end{paso}

\begin{paso}{4: Tabla de signos de $q(x) = x^2 - x - 1$}
\begin{center}
\renewcommand{\arraystretch}{1.6}
\begin{tabular}{c|ccccc}
    & $x{<}r_1$ & $x{=}r_1$ & $r_1{<}x{<}r_2$ & $x{=}r_2$ & $x{>}r_2$ \\
\hline
$(x-r_1)$ & $-$ & $0$ & $+$ & $+$ & $+$ \\
$(x-r_2)$ & $-$ & $-$ & $-$ & $0$ & $+$ \\
\hline
\rowcolor{celesteSuave}
$q(x)$ & $+$ & $0$ & $\mathbf{-}$ & $0$ & $+$ \\
\end{tabular}
\end{center}
\[
    q(x) < 0 \quad\Longleftrightarrow\quad
    x \in \left(\frac{1-\sqrt{5}}{2},\;\frac{1+\sqrt{5}}{2}\right)
\]

\textbf{Recta numérica:}

\begin{center}
\begin{tikzpicture}[scale=1.1]
    \draw[->] (-1.5,0) -- (3.5,0) node[right] {$x$};
    \fill[azulFuerte] (-0.618,0) circle (2.2pt)
        node[below=4pt] {\small$r_1{\approx}{-}0.618$};
    \fill[azulFuerte]  (1.618,0) circle (2.2pt)
        node[below=4pt] {\small$r_2{\approx}1.618$};
    \node at (-1.0, 0.35) {\color{rojoAlerta}\small$q>0$};
    \node at ( 0.5, 0.35) {\color{verdePaso}\small$q<0$};
    \node at ( 2.6, 0.35) {\color{rojoAlerta}\small$q>0$};
    \draw[verdePaso, very thick] (-0.618,0) -- (1.618,0);
    \draw[azulFuerte!60, dashed] (1,-0.5) -- (1,0.5);
    \draw[azulFuerte!60, dashed] (2,-0.5) -- (2,0.5);
    \fill[azulFuerte] (1,0) circle (2.2pt);
    \fill[azulFuerte!60] (2,0) circle (2.2pt);
    \node[azulFuerte] at (1,-0.55) {\small$1$};
    \node[azulFuerte] at (2,-0.55) {\small$2$};
    \node[azulFuerte] at (1.5,-0.9) {\small Dominio $[\,1,2\,]$};
\end{tikzpicture}
\end{center}
\end{paso}

\begin{paso}{5: Intersección con el dominio}
\[
    \left(\frac{1-\sqrt{5}}{2},\;\frac{1+\sqrt{5}}{2}\right)
    \;\bigcap\;[\,1,\,2\,]
    \;=\;
    \left[\,1,\;\frac{1+\sqrt{5}}{2}\right)
\]
Verificación: $q(1)=1-1-1=-1<0$ \checkmark;\quad
$q\!\left(\tfrac{1+\sqrt5}{2}\right)=0$, no cumple la desigualdad estricta \checkmark.
\end{paso}

\resultado{x \;\in\; \left[\,1,\;\dfrac{1+\sqrt{5}}{2}\right) \;\approx\; [\,1,\;1.618\,)}

\clearpage

% ─── EJERCICIO 2.4 ──────────────────────────────────────────
\exheader{2.4}{\dfrac{1}{\sqrt{3 - \sqrt{x+1}}} > 1}

\begin{paso}{1: Dominio natural}
\begin{itemize}[itemsep=0.2em]
    \item $x + 1 \geq 0 \;\Rightarrow\; x \geq -1$
    \item $3 - \sqrt{x+1} > 0 \;\Rightarrow\; \sqrt{x+1} < 3 \;\Rightarrow\; x+1 < 9 \;\Rightarrow\; x < 8$
\end{itemize}
Dominio: $x \in [-1,\,8)$.
\end{paso}

\begin{paso}{2: Primer paso algebraico}
Como $\sqrt{3-\sqrt{x+1}} > 0$, la fracción $> 1$ equivale a que el denominador sea $< 1$:
\[
    \frac{1}{\sqrt{3-\sqrt{x+1}}} > 1
    \quad\Longleftrightarrow\quad
    \sqrt{3-\sqrt{x+1}} < 1
\]
\end{paso}

\begin{paso}{3: Primera elevación al cuadrado (ambos lados $\geq 0$)}
\[
    3 - \sqrt{x+1} < 1
    \quad\Longrightarrow\quad
    \sqrt{x+1} > 2
\]
\end{paso}

\begin{paso}{4: Segunda elevación al cuadrado}
\[
    \sqrt{x+1} > 2
    \quad\stackrel{(\,)^2}{\Longleftrightarrow}\quad
    x + 1 > 4
    \quad\Longrightarrow\quad
    x > 3
\]
\end{paso}

\begin{paso}{5: Intersección con el dominio}
\[
    [-1,\,8) \;\cap\; (3,+\infty) \;=\; (3,\,8)
\]
\end{paso}

\resultado{x \;\in\; (3,\;8)}

\clearpage

% ─── EJERCICIO 2.5 ──────────────────────────────────────────
\exheader{2.5}{\dfrac{\sqrt{x - \sqrt{x}}}{x - 3} \geq 0}

\begin{paso}{1: Dominio natural}
\textbf{Condición 1:} $x \geq 0$ (para que $\sqrt{x}$ exista).\\[0.3em]
\textbf{Condición 2:} $x - \sqrt{x} \geq 0$. Sea $u = \sqrt{x} \geq 0$:
\[
    u^2 - u \geq 0
    \quad\Longleftrightarrow\quad
    u(u-1) \geq 0
    \quad\Longleftrightarrow\quad
    u = 0 \;\text{ ó }\; u \geq 1
    \quad\Longleftrightarrow\quad
    x = 0 \;\text{ ó }\; x \geq 1
\]
\textbf{Condición 3:} $x \neq 3$ (denominador $\neq 0$).\\[0.3em]
Dominio: $\{0\} \cup [\,1,+\infty) \setminus \{3\}$.
\end{paso}

\begin{paso}{2: Análisis de la fracción}
El numerador $\sqrt{x-\sqrt{x}} \geq 0$ \textbf{siempre} en el dominio.
Es $= 0$ cuando $x - \sqrt{x} = 0$, es decir, $x = 0$ ó $x = 1$.

\begin{caso}{Caso A --- Fracción $= 0$: numerador nulo}
$x = 0$: denominador $= -3 \neq 0$. Fracción $= 0 \geq 0$. \textbf{Válido.}\\
$x = 1$: denominador $= -2 \neq 0$. Fracción $= 0 \geq 0$. \textbf{Válido.}
\end{caso}

\bigskip

\begin{caso}{Caso B --- Fracción $> 0$: ambos factores positivos}
Numerador $> 0 \;\Rightarrow\; x > 1$ (en el dominio).\\
Denominador $> 0 \;\Rightarrow\; x > 3$.\\
Combinando: $x > 3$.\\[0.2em]
\textbf{Solución B:} $(3,+\infty)$.
\end{caso}

\end{paso}

\begin{paso}{3: Unión de casos}
\[
    \{0,\,1\} \;\cup\; (3,\,+\infty)
\]
\end{paso}

\resultado{x \;\in\; \{0,\,1\}\;\cup\;(3,\,+\infty)}

\clearpage

% ============================================================
%  SECCIÓN III — RAÍZ DENTRO DE RAÍZ
% ============================================================
\secheader{SECCIÓN III --- Raíz dentro de Raíz}

% ─── EJERCICIO 3.1 ──────────────────────────────────────────
\exheader{3.1}{\sqrt{\sqrt{x} - 1} \leq 2}

\begin{paso}{1: Dominio natural}
$\sqrt{x} - 1 \geq 0 \;\Rightarrow\; \sqrt{x} \geq 1 \;\Rightarrow\; x \geq 1$.\\
Dominio: $[\,1,+\infty)$.
\end{paso}

\begin{paso}{2: Primera elevación al cuadrado (LHS $\geq 0$, RHS $= 2 > 0$)}
\[
    \sqrt{\sqrt{x}-1} \leq 2
    \quad\stackrel{(\,)^2}{\Longleftrightarrow}\quad
    \sqrt{x} - 1 \leq 4
    \quad\Longrightarrow\quad
    \sqrt{x} \leq 5
\]
\end{paso}

\begin{paso}{3: Segunda elevación al cuadrado}
\[
    \sqrt{x} \leq 5
    \quad\stackrel{(\,)^2}{\Longleftrightarrow}\quad
    x \leq 25
\]
\end{paso}

\begin{paso}{4: Intersección con el dominio}
\[
    [\,1,+\infty) \;\cap\; (-\infty,\,25\,] \;=\; [\,1,\,25\,]
\]
\end{paso}

\resultado{x \;\in\; [\,1,\;25\,]}

\bigskip

% ─── EJERCICIO 3.2 ──────────────────────────────────────────
\exheader{3.2}{\sqrt{2 + \sqrt{x-1}} < 2}

\begin{paso}{1: Dominio natural}
$x - 1 \geq 0 \;\Rightarrow\; x \geq 1$. La expresión $2+\sqrt{x-1} \geq 0$ siempre.\\
Dominio: $[\,1,+\infty)$.
\end{paso}

\begin{paso}{2: Primera elevación al cuadrado (ambos lados $\geq 0$)}
\[
    \sqrt{2+\sqrt{x-1}} < 2
    \quad\stackrel{(\,)^2}{\Longleftrightarrow}\quad
    2 + \sqrt{x-1} < 4
    \quad\Longrightarrow\quad
    \sqrt{x-1} < 2
\]
\end{paso}

\begin{paso}{3: Segunda elevación al cuadrado}
\[
    \sqrt{x-1} < 2
    \quad\stackrel{(\,)^2}{\Longleftrightarrow}\quad
    x - 1 < 4
    \quad\Longrightarrow\quad
    x < 5
\]
\end{paso}

\begin{paso}{4: Intersección con el dominio}
\[
    [\,1,+\infty) \;\cap\; (-\infty,\,5) \;=\; [\,1,\,5)
\]
\end{paso}

\resultado{x \;\in\; [\,1,\;5)}

\clearpage

% ─── EJERCICIO 3.3 ──────────────────────────────────────────
\exheader{3.3}{\sqrt{x + \sqrt{x-1}} \geq 2}

\begin{paso}{1: Dominio natural}
$x - 1 \geq 0 \;\Rightarrow\; x \geq 1$. La expresión $x + \sqrt{x-1} \geq 0$ siempre.\\
Dominio: $[\,1,+\infty)$.
\end{paso}

\begin{paso}{2: Primera elevación al cuadrado (ambos lados $\geq 0$)}
\[
    \sqrt{x+\sqrt{x-1}} \geq 2
    \quad\stackrel{(\,)^2}{\Longleftrightarrow}\quad
    x + \sqrt{x-1} \geq 4
    \quad\Longrightarrow\quad
    \sqrt{x-1} \geq 4 - x
\]
\end{paso}

\begin{paso}{3: Análisis por casos}

\begin{caso}{Caso A --- $4 - x \leq 0$, es decir $x \geq 4$}
$\sqrt{x-1} \geq 0 \geq 4-x$ automáticamente.\\
\textbf{Solución A:} $[\,4,+\infty)$.
\end{caso}

\bigskip

\begin{caso}{Caso B --- $4 - x > 0$, es decir $1 \leq x < 4$}
Ambos lados $\geq 0$, elevamos al cuadrado:
\[
    x - 1 \geq (4-x)^2 = 16 - 8x + x^2
    \quad\Longrightarrow\quad
    x^2 - 9x + 17 \leq 0
\]
Discriminante: $\Delta = 81 - 68 = 13$.
\[
    r_{1,2} = \frac{9 \pm \sqrt{13}}{2}, \qquad
    r_1 = \frac{9-\sqrt{13}}{2} \approx 2.697, \quad
    r_2 = \frac{9+\sqrt{13}}{2} \approx 6.303
\]

\begin{center}
\renewcommand{\arraystretch}{1.5}
\begin{tabular}{c|ccccc}
    & $x{<}r_1$ & $x{=}r_1$ & $r_1{<}x{<}r_2$ & $x{=}r_2$ & $x{>}r_2$ \\
\hline
\rowcolor{celesteSuave}
$x^2-9x+17$ & $+$ & $0$ & $\mathbf{-}$ & $0$ & $+$ \\
\end{tabular}
\end{center}

$x^2-9x+17 \leq 0 \;\Leftrightarrow\; x \in [\,r_1, r_2\,]$.
Intersectado con $[\,1,\,4)$ (nótese $r_1 \approx 2.697 \in [\,1,4)$ y $r_2 \approx 6.303 > 4$):

\textbf{Solución B:} $x \in \left[\dfrac{9-\sqrt{13}}{2},\;4\right)$.
\end{caso}

\end{paso}

\begin{paso}{4: Unión de soluciones}
\[
    \left[\frac{9-\sqrt{13}}{2},\;4\right) \;\cup\; [\,4,+\infty)
    \;=\;
    \left[\frac{9-\sqrt{13}}{2},\;+\infty\right)
\]
\end{paso}

\resultado{x \;\in\; \left[\dfrac{9-\sqrt{13}}{2},\;+\infty\right) \;\approx\; [\,2.697,\;+\infty)}

\clearpage

% ─── EJERCICIO 3.4 ──────────────────────────────────────────
\exheader{3.4}{\sqrt{3 - \sqrt{2+x}} \geq 1}

\begin{paso}{1: Dominio natural}
\begin{itemize}[itemsep=0.2em]
    \item $2 + x \geq 0 \;\Rightarrow\; x \geq -2$
    \item $3 - \sqrt{2+x} \geq 0 \;\Rightarrow\; \sqrt{2+x} \leq 3
          \;\Rightarrow\; 2+x \leq 9 \;\Rightarrow\; x \leq 7$
\end{itemize}
Dominio: $[\,-2,\,7\,]$.
\end{paso}

\begin{paso}{2: Elevación al cuadrado (ambos lados $\geq 0$ en el dominio)}
\[
    \sqrt{3-\sqrt{2+x}} \geq 1
    \quad\stackrel{(\,)^2}{\Longleftrightarrow}\quad
    3 - \sqrt{2+x} \geq 1
    \quad\Longrightarrow\quad
    \sqrt{2+x} \leq 2
\]
\end{paso}

\begin{paso}{3: Segunda elevación al cuadrado}
\[
    \sqrt{2+x} \leq 2
    \quad\stackrel{(\,)^2}{\Longleftrightarrow}\quad
    2 + x \leq 4
    \quad\Longrightarrow\quad
    x \leq 2
\]
\end{paso}

\begin{paso}{4: Intersección con el dominio}
\[
    [\,-2,\,7\,] \;\cap\; (-\infty,\,2\,] \;=\; [\,-2,\,2\,]
\]
\end{paso}

\resultado{x \;\in\; [\,-2,\;2\,]}

\bigskip

% ─── EJERCICIO 3.5 ──────────────────────────────────────────
\exheader{3.5}{\sqrt{1 - \sqrt{1-x^2}} \leq \dfrac{\sqrt{2}}{2}}

\begin{paso}{1: Dominio natural}
\begin{itemize}[itemsep=0.2em]
    \item $1 - x^2 \geq 0 \;\Rightarrow\; x^2 \leq 1 \;\Rightarrow\; -1 \leq x \leq 1$
    \item $1 - \sqrt{1-x^2} \geq 0 \;\Rightarrow\; \sqrt{1-x^2} \leq 1$, siempre cierto.
\end{itemize}
Dominio: $[\,-1,\,1\,]$.
\end{paso}

\begin{paso}{2: Primera elevación al cuadrado (ambos lados $\geq 0$)}
\[
    \sqrt{1-\sqrt{1-x^2}} \leq \frac{\sqrt{2}}{2}
    \quad\stackrel{(\,)^2}{\Longleftrightarrow}\quad
    1 - \sqrt{1-x^2} \leq \frac{1}{2}
    \quad\Longrightarrow\quad
    \sqrt{1-x^2} \geq \frac{1}{2}
\]
\end{paso}

\begin{paso}{3: Segunda elevación al cuadrado}
\[
    \sqrt{1-x^2} \geq \frac{1}{2}
    \quad\stackrel{(\,)^2}{\Longleftrightarrow}\quad
    1 - x^2 \geq \frac{1}{4}
    \quad\Longrightarrow\quad
    x^2 \leq \frac{3}{4}
    \quad\Longrightarrow\quad
    |x| \leq \frac{\sqrt{3}}{2}
\]
Es decir, $x \in \left[-\dfrac{\sqrt{3}}{2},\;\dfrac{\sqrt{3}}{2}\right]$.
\end{paso}

\begin{paso}{4: Intersección con el dominio}
$\left[-\dfrac{\sqrt{3}}{2},\;\dfrac{\sqrt{3}}{2}\right] \subset [\,-1,\,1\,]$, luego la intersección no modifica el intervalo.
\end{paso}

\resultado{x \;\in\; \left[-\dfrac{\sqrt{3}}{2},\;\dfrac{\sqrt{3}}{2}\right] \;\approx\; [\,-0.866,\;0.866\,]}

\vspace{1.2cm}
\begin{center}
    \small\textbf{Usach Premium} --- ¡Nos vemos en la ayudantía! \\
    \textit{"Por y para estudiantes."}
\end{center}

\end{document}
